\subsection{The integral Fermat comparison and its actual lift functor}

We distinguish the full signed automorphism input of \S49 from its
explicit \emph{minimal-stage} specialization in \S51.2.  In the latter,
the conjugators lie in the rank-two vector group, with the reversing
component recorded separately.  The precise formulas used below are
\S51.2, folios 665 verso--669, C4 scans 36--43, equations (15),
(18), (21), (28), (30), (31), and (38)
\citep{Grothendieck1981FermatStage}.  The supplied image of scan 41
was inspected directly.  Its doubtful signs do not affect the
characteristic-two calculations, and our counterexample has all sign
parameters zero.  Retaining a conjugator sign is not the same as
retaining every infinitesimal automorphism.

\begin{statement}{B.24a}{Proved}{The canonical integral projection has a finite image at two}
Fix the discrete modular group $C=\mathrm{SL}_2(\Z)$ with
\[
 \rho=\begin{pmatrix}0&1\\-1&1\end{pmatrix},\qquad
 \sigma=\begin{pmatrix}0&-1\\1&0\end{pmatrix},\qquad
 \omega=\rho^3=\sigma^2=-I.
\]
Put $x=\omega(\sigma\rho)^2$, $y=\omega(\rho\sigma)^2$,
$F=\langle x,y\rangle$, $K=\ker(F\to(\Z/2)^2)$, and
$M=F^{\mathrm{ab}}$.  Then $F$ is free of rank two and index
twelve in $C$, while $K=\Gamma(4)$ is free of rank five and
index forty-eight.  Applying the integral cocycle construction of
B.20 to $C/[F,F]$ and $C/[K,K]$ gives group schemes
\[
\begin{split}
 1&\longrightarrow\Ga^2\longrightarrow\mathscr C_0
       \longrightarrow H_{12,\Z}\longrightarrow1,\\
 1&\longrightarrow\Ga^5\longrightarrow\mathscr C_2
       \longrightarrow H_{48,\Z}\longrightarrow1.
\end{split}
\]
Here $H_{12}=C/F$ and $H_{48}=C/K$.
There is a canonical homomorphism of integral group schemes
$\pi:\mathscr C_2\to\mathscr C_0$ induced by the identity of $C$.
The inverse image of $(F/K)_\Z\subset H_{48,\Z}$ is the integral
abelian-kernel stage for $q_2:F\to(\Z/2)^2$.  Its rational fibre is
the quotient of $\mathscr R_{(\Z/2)^2}(F)$ obtained by abelianizing
the unipotent radical.

In the Schreier basis
\[
\begin{aligned}
 k_1&=x^2,& k_2&=yxy^{-1}x^{-1},\\
 k_3&=y^2,& k_4&=xyxy^{-1},\\
 k_5&=xy^2x^{-1}.
\end{aligned}
\]
the map $j:K^{\mathrm{ab}}\to M$ is
\[
 J=\begin{pmatrix}2&0&0&2&0\\0&0&2&0&2\end{pmatrix}.
 \tag{B.24a.1}
\]
Consequently, over $\F_2$, $\pi$ factors as
\[
 \mathscr C_{2,\F_2}\longrightarrow H_{48,\F_2}
       \xrightarrow{\ i\ }\mathscr C_{0,\F_2},
 \tag{B.24a.2}
\]
where the first map is faithfully flat and $i$ is a closed immersion.
In the integral vector coordinates on each of the twelve components of
$\mathscr C_0$, the image $Y=i(H_{48,\F_2})$ is cut out by
$X^2-X=Y_1^2-Y_1=0$.  This is the scheme-theoretic image, over every
$\F_2$-algebra, not merely a statement about rational points.
\end{statement}
\begin{proof}
The modular presentation in E.2 gives
$C=C_6*_{C_2}C_4$.  Direct multiplication shows that conjugation
by $\rho,\sigma$ acts on $x,y$ by
\[
 \rho:(x,y)\mapsto(y,x^{-1}y^{-1}),\qquad
 \sigma:(x,y)\mapsto(y,x).
 \tag{B.24a.3}
\]
Thus $F$ is normal.  Killing $x,y$ gives the presentation
$\langle r,s\mid r^6=1, s^2=r^3, srs^{-1}=r^{-1}\rangle$.
Every word has the form $r^is^e$, $0\le i<6$, $0\le e<2$.
Reduction modulo four realizes all twelve forms: the images of $x,y$
generate a normal Klein four subgroup of the forty-eight matrices in
$\mathrm{SL}_2(\Z/4)$.  Hence the index is twelve.  The two finite
amalgam factors inject into this quotient.  The subgroup therefore acts
freely on the amalgam tree; the quotient has five vertices and six
edges, and its fundamental group has rank two.  Its two indicated
generators are a free basis, also obtained directly by Schreier rewriting.
Their distinct images generate that Klein four group.  It follows that
the kernel of $F\to(\Z/2)^2$ is exactly the kernel of reduction
$C\to\mathrm{SL}_2(\Z/4)$, and Schreier gives rank five.

Take the four Schreier representatives $1,x,y,xy$ in $F$.
The five nontrivial edges give exactly the words $k_i$ in the
statement.  Their exponent sums prove (B.24a.1).  In particular
$\operatorname{im}j=2\Z^2$ and its kernel is free of rank three.
This argument determines the integral lattice, rather than only
the rational rank.

Here is the precise scheme construction and comparison.  For either
$L=F^{\mathrm{ab}}$ or $L=K^{\mathrm{ab}}$, choose actual word
representatives $s_h\in C$ of the relevant finite quotient $H$.
Let $A_h$ be conjugation on $L$ and let
$c(h,k)$ be the class of $s_hs_ks_{hk}^{-1}$.  On
$\coprod_{h\in H}\Ga^{\operatorname{rank}L}$ set
\[
 (v,h)(w,k)=(v+A_hw+c(h,k),hk).
 \tag{B.24a.4}
\]
The word identity for $c$ proves associativity, including the integral
coefficients.  Different representatives translate the coordinates and
give the unique compatible isomorphism.  Write $p:H_{48}\to H_{12}$
and let $m_h\in M$ be the class of $s_hs_{p(h)}^{-1}$.  Then
\[
                    \pi(v,h)=(Jv+m_h,p(h)).
 \tag{B.24a.5}
\]
Equivariance of $j$ and the identity
\[
 Jc_2(h,k)+m_{hk}
   =m_h+A_{p(h)}m_k+c_0(p(h),p(k))
\]
prove that this is a group-scheme homomorphism.  The construction is
canonical because it is induced by the homomorphism of the two discrete
extensions, with the lattice maps extended linearly.  Restriction to the
four components above $F/K$ is the same construction for $F/[K,K]$.
Over $\Q$ the universal property in B.8 identifies it with the indicated
abelian-radical quotient of relative completion.

Modulo two the linear part $J$ vanishes.  The kernel of
$C\to\mathscr C_0(\F_2)$ consists precisely of the words in $F$
whose two exponent sums are even, namely $K$.  Its image is therefore
$H_{48}$.  On each component of $\mathscr C_0$ this is exactly the
four-vector grid $\F_2^2$.  The grid is the closed constant scheme
defined by $X^2-X,Y_1^2-Y_1$; these two polynomials are separable.
This proves the closed immersion and factorization.  Projection from
the forty-eight vector-group components is faithfully flat, so its
composite has precisely this scheme-theoretic image.

The program \texttt{g\_signed\_fermat2.py} constructs the forty-eight
word representatives, the integral action matrices and the
entire cocycle.  It certifies each Schreier word by its exact integral
$2\times2$ matrix, checks all $48^3=110592$ cocycle identities and all
$48^2=2304$ projection identities, and records the complete model in
\texttt{g\_signed\_fermat2\_model.json}.  Thus the finite image assertion
is checked on the actual integral modular extension.
\end{proof}

\begin{statement}{B.24b}{Proved}{The liftable locus and the marked lift kernel at two}
Let $k=\F_2$.  Retain the two \emph{constant} sign coordinates
$\nu_2,\nu_3\in(\Z/2)_k$ of the fine parameter scheme in C4 scan 41.
Its $k$-fibre is the disjoint union of four copies of $\Ga^3$, with
coordinates $a,b,k_0$.  Its explicitly linear action on $\mathscr C_0$
factors through the group scheme
\[
 \mathscr U=\Ga\times(\Z/2)_k,\qquad
 (t,e)=(a+\nu_2+\nu_3,\nu_2),
 \tag{B.24b.1}
\]
where
\[
\begin{split}
 u_{t,e}(\rho)&=(t,t+e)\rho,\\
 u_{t,e}(\sigma)&=(t,t)\omega^e\sigma,\\
 u_{t,e}|_{M\otimes\Ga}&=1.
\end{split}
\tag{B.24b.2}
\]
The fibres of the fine parameter map are
$\Ga^2\times(\Z/2)_k$; the last factor retains $\nu_3$.

The functor of pairs $(u,\Phi)$ with $u\in\mathscr U$ and
$\Phi\in\Aut(\mathscr C_2)$ satisfying
$\pi\Phi=u\pi$ has image exactly
\[
       \mathscr U_{\mathrm{lift}}
       =V(t^2-t)\times(\Z/2)_k\simeq(\Z/2)_k^2.
 \tag{B.24b.3}
\]
This assertion holds on every $k$-algebra, and also for fppf-local lifts.
Every nonempty lift fibre is a torsor under
$\Aut_{H_{48,k}}(\mathscr C_{2,k})$.

There is a completely explicit marked version.  Label the six boundary
vector subgroups by the classes of
\[
 x^2,\quad yx^2y^{-1},\quad y^2,\quad xy^2x^{-1},
 \quad z^2,\quad xz^2x^{-1},\qquad z=(xy)^{-1}.
 \tag{B.24b.4}
\]
Require a lift of $(t,e)\in\mathscr U_{\mathrm{lift}}$ to induce on
these six labels the permutation of
$\operatorname{Int}(x)^t\tau^e$, where
$\tau=\operatorname{Int}(\operatorname{diag}(-1,1))$.
The kernel of this marked lift functor is the ind-affine group functor
\[
 \mathscr K_{\mathrm{mark}}\simeq\Ga^6\rtimes\mathscr A_1,
 \qquad
 \mathscr A_1(R)=
 \{f\in\Aut_{R\text{-groups}}(\mathbb G_{a,R}):f(1)=1\}.
 \tag{B.24b.5}
\]
The action of $f$ on the six affine coordinates is coordinatewise
application of its additive polynomial.  Each marked lift fibre is a
torsor under this kernel, and the whole marked lift functor is
represented by four copies of this ind-affine functor.

For $A=k[\epsilon]/\epsilon^2$, the fine scheme has $256$ points,
its minimal automorphism image has $8$, and exactly $4$ of those
automorphisms lift.  Thus precisely $128$ fine parameter points lift.
At the identity, the kernel of degree-at-most-two polynomial lifts
reducing to the identity has $k$-dimension seven.  Its one-dimensional
subspace fixing every discrete word is
\[
        (v,h)\longmapsto
           \bigl(v+\epsilon a(v^{[2]}-v),h\bigr),\qquad a\in k.
 \tag{B.24b.6}
\]
Allowing arbitrary finite additive-polynomial degree makes this last
infinitesimal subspace $\bigoplus_{r\ge1}k$, with basis
$v^{[2^r]}-v$.
\end{statement}
\begin{proof}
In the basis $(x,y)$, conjugation has matrices
\[
 R=\begin{pmatrix}0&-1\\1&-1\end{pmatrix},\qquad
 S=\begin{pmatrix}0&1\\1&0\end{pmatrix}.
\]
The coherent integral scan-41 equations are
\[
 \nu=2\nu'+\nu_2,\qquad
 (1-R)\alpha_0-\nu_3(0,1)^t
   =(1-S)\beta_0+\nu(0,1)^t,\qquad \mu=1+2\nu.
\]
Writing $d=\nu_2+\nu_3$, their integral solution is
\[
 \alpha_0=(a,d+2k_0),\qquad
 \beta_0=(a+d+2k_0+b,b),\qquad \nu'=d+3k_0.
\]
The unit condition is $12k_0+6\nu_2+4\nu_3+1\ne0$ in the
sense of invertibility.  In characteristic two it is automatically
one.  Substitution gives (B.24b.1)--(B.24b.2); the same equations with
the alternative coherent signs give the identical characteristic-two
formulas.  They define actual automorphisms: over $k$,
\[
                    u_{t,e}=\operatorname{Int}(t x)\tau^e
 \tag{B.24b.7}
\]
on $\mathscr C_0$.  Here $t x$ is a vector-group element, not a
formal fractional discrete word.  The matrix $\tau$ acts on $M$ as
$-1$, hence as the identity in characteristic two; it sends
$\rho\mapsto(-y)\rho$, $\sigma\mapsto\omega\sigma$.
Consequently (B.24b.7) has the direct-product law claimed for
$\mathscr U$.  Conversely the image of $\rho$ determines $t$
and that of $\sigma$ determines $e$.  Solving for
$a=t+\nu_2+\nu_3$ leaves $b,k_0,\nu_3$ arbitrary, proving the
fine-parameter fibre assertion.

Any compatible lift must preserve the scheme-theoretic image $Y$ of
B.24a.  The first coordinate of $u_{t,e}(\rho)$ is $t$, so the equation
$X^2-X=0$ on $Y$ forces $t^2=t$.  A faithfully flat ring extension
cannot make a nonzero $t^2-t$ zero.  Conversely an idempotent $t$
is locally the constant zero or one.  On these open-and-closed pieces
$\operatorname{Int}(x)^t\tau^e$ is an actual lift on $\mathscr C_2$:
both maps come from automorphisms of $C$ preserving $K$, and their
lattice actions extend integrally.  This proves (B.24b.3), with no
assumption about reduced bases.  If two lifts of the same $u$ are given,
their ratio lies over the identity of $\mathscr C_0$.  Since the
factor $i$ in (B.24a.2) is a monomorphism, this says exactly that the
ratio induces the identity on $H_{48}$.  Composition gives the stated
torsor action.

We compute that kernel with the extra six markings.  Write
$W=K^{\mathrm{ab}}\otimes k$.  Every automorphism over $H_{48}$
has the unique form
\[
             (v,h)\longmapsto(T(v)+b_h,h),\qquad b_1=0,
 \tag{B.24b.8}
\]
where $T$ is an additive-polynomial automorphism of $W\otimes\Ga$.
Equation (B.24a.4) is respected if and only if
\[
 T A_h=A_hT,\qquad
 b_h+A_hb_k-b_{hk}=T(c(h,k))-c(h,k).
 \tag{B.24b.9}
\]
These are necessary by multiplying general points, and sufficient by
the same identity.  They are an explicit description of the whole
unmarked kernel as well, with its nonsemisimple deck action retained.

The six vectors $v_i$ of (B.24b.4), in the basis $k_i$, are
\[
\begin{gathered}
 (1,0,0,0,0),\quad(0,1,0,1,0),\quad(0,0,1,0,0),\\
 (0,0,0,0,1),\quad(0,0,-1,-1,0),\quad(-1,-1,0,0,-1).
\end{gathered}
\]
Their sole relation is $\sum_i v_i=0$, and any five are an integral
basis: the six complementary determinants alternate between $-1$ and
$1$.  Additivity of a polynomial over a $k$-algebra eliminates every
mixed-variable monomial.  For a monomial in one variable, a nonzero
intermediate binomial coefficient modulo two eliminates its coefficient
unless its degree is a power of two.  Thus, over an arbitrary
$k$-algebra, we can write
$T(v)=\sum_r D_r v^{[2^r]}$.  Preservation of each labelled line
forces $D_rv_i=a_{r,i}v_i$.  Apply $D_r$ to the sole relation.
Its unimodularity forces all six $a_{r,i}$ to be equal, even over a
nonreduced ring.  Thus $D_r=a_rI$ and $T$ is coordinatewise application
of one additive polynomial $f$.  This also proves equivariance,
since every action matrix has entries in $k$.

Put $\lambda=f(1)$.  The cocycle entries belong to the prime field,
so (B.24b.9) becomes
$\delta b=(\lambda-1)c$.  In fact $\lambda=1$ is necessary.
Let $X$ be conjugation by the image of $x$ in $H_{48}$, which has
order two, and choose its representative to be $x$.
Then $c(x,x)=k_1$.  The row vector
\[
                    \ell=(1,0,0,1,0)\quad\text{over }k
\]
satisfies $\ell(1+X)=0$ and $\ell(k_1)=1$.
Apply $\ell$ to $\delta b(x,x)=(\lambda-1)k_1$ to obtain
$\lambda-1=0$.  Changing the representative adds a norm to
$c(x,x)$ and does not change this argument.  The exact matrix $X$
and this norm certificate are in the output.  Therefore the possible
translations are precisely $Z^1(H_{48},W\otimes R)$, independently
of $f\in\mathscr A_1(R)$.

This cocycle functor is $\Ga^6$.  Here is a small reproducible rank
certificate, avoiding a semisimplicity assumption at two.  A cocycle
is determined by $d_\rho,d_\sigma\in R^5$.  Propagate these two
values along the forty-eight fixed representative words and impose
the two generator edges at each vertex of the Cayley graph.  The
resulting linear equations have the following reduced row basis on
$(d_{\rho,1},\ldots,d_{\rho,5},d_{\sigma,1},\ldots,d_{\sigma,5})$:
\[
 \begin{pmatrix}
 1&0&0&1&0&0&0&1&0&1\\
 0&1&1&1&1&0&0&1&0&1\\
 0&0&0&0&0&1&0&1&0&0\\
 0&0&0&0&0&0&0&0&1&1
 \end{pmatrix}.
 \tag{B.24b.10}
\]
The edge identities are sufficient: propagation along any word now
gives the prescribed value at its finite group element, hence gives
the cocycle identity for all products.  Four unit pivots prove the
claim over every $k$-algebra.  The action matrices also have no common
fixed vector.  Thus coboundaries have dimension five and
$\dim_kH^1(H_{48},W)=1$; this does not enter as an uncomputed
lifting assumption.  Composition in (B.24b.8) is
$(b,f)(b',f')=(b+f(b'),f\circ f')$.  A basis of the cocycle space
over $k$ identifies this with (B.24b.5).

For clarity, representability here is \emph{ind-affine}, not finite
type.  Bound the degrees of both an additive polynomial and its
polynomial inverse; their coefficients satisfy finitely many polynomial
equations for $f\circ g=g\circ f=T$ and $f(1)=1$.  These are affine
parameter schemes.  Increasing the degree bounds by adjoining zero
coefficients gives the required ind-affine group functor.  Formula
(B.24b.8) and the finite set of equations (B.24b.9) give the same
description for the unmarked kernel.  The four chosen lifts above
trivialize the marked torsors as schemes on the four constant components;
they are not claimed to form a homomorphic section.

Over the dual numbers, every additive polynomial reducing to the
identity and fixing one has the unique form
\[
             f(T)=T+\epsilon\sum_{r\ge1}a_r(T^{2^r}-T),
             \qquad a_r\in k,
\]
with finite support.  Its inverse is obtained by changing the sign
of the correction, which is the same sign in characteristic two.
The translation cocycles reducing to zero contribute
$\epsilon k^6$.  Restricting to degree at most two leaves one
additional coefficient, proving dimension seven.  If every discrete
word is fixed, the representative values force $b_h=0$ and the
integral basis values force $f(1)=1$; this is precisely the displayed
Frobenius family.  The program checks additivity, both inverse identities,
and the nonzero $\epsilon T^2$ coefficient as polynomial identities,
not just on the four $A$-points.  It enumerates all $256$ fine parameter
points: each of the eight minimal automorphisms has $32$ preimages,
and $t^2=t$ leaves four automorphisms and $128$ parameter points.

Finally, $\nu_3$ is a locally constant label.  Over these dual numbers
$(\Z/2)(A)$ has two elements, while $\mu_2(A)=\{1,1+\epsilon\}$.
The map $\nu_3\mapsto(-1)^{\nu_3}$ sends both constant labels to $1$.
Replacing the constant sign scheme by $\mu_2$ would therefore change
the functor.  Neither that replacement nor retaining $\nu_3$ repairs
the obstruction $t^2-t\ne0$.
\end{proof}

\begin{statement}{B.24c}{Proved}{The rational comparison and the direction of functoriality}
For the inner subgroup $u_t=\operatorname{Int}(t x)$ of the minimal
model, a compatible homomorphic lift exists over $\Z[1/2]$:
\[
 \mathbb G_{a,\Z[1/2]}\longrightarrow\Aut(\mathscr C_2)_{\Z[1/2]},
 \qquad t\longmapsto\operatorname{Int}\bigl((t/2)k_1\bigr).
 \tag{B.24c.1}
\]
Over $\Q$, its inner lifts supplied by elements of the identity vector
subgroup, at a fixed $t$, form a $\Ga^3$-torsor.
Every minimal fine datum over a $\Q$-algebra has a compatible
automorphism lift at this first abelian relative stage, after fixed
rational splitting choices.  This existence does not make the lift
unique or extend it to the integral base at two.

Separately, an actual automorphism of the discrete group $C$
preserving $F$ and $K$ induces compatible automorphisms of both
integral stages.  An actual continuous profinite automorphism preserving
the corresponding characteristic quotients acts on their finite/profinite
quotients.  For the rational relative completion, functoriality applies
to an actual homomorphism of its source with a compatible automorphism
of its finite reductive quotient.  An automorphism of a coarser completed
quotient is not such input.
\end{statement}
\begin{proof}
Equation (B.24a.1) sends $(t/2)k_1$ to $t x$, proving compatibility
in (B.24c.1).  The vector subgroup is commutative, so this is a group
homomorphism.  Over $\Q$, all preimages of $t x$ in the identity
vector subgroup form a torsor under
$\ker(j\otimes\Q)\simeq\Ga^3$.  The associated inner automorphisms
are distinct.  Indeed $W_\Q^{H_{48}}=0$: the action permutes the six
boundary vectors transitively, and $W_\Q$ is the quotient of that
permutation module by its constant line.  Averaging over the finite
group proves that this quotient has no invariant vector.  Equivalently,
the exact action matrices in the output have common fixed-space rank
zero over $\Q$.

We give the promised rational existence argument without calling it
canonical.  A finite-group vector extension over a $\Q$-algebra splits:
for its cocycle $c$, put
$b(h)=|H|^{-1}\sum_k c(h,k)$.  The cocycle identity gives
$c(h,k)=b(h)+A_hb(k)-b(hk)$.  This constructs the splitting with
rational coefficients.  Splittings of $\mathscr C_2$ and
$\mathscr C_0$ can be made compatible with $\pi$.  The discrepancy
of two chosen splittings is a one-cocycle in $M_\Q$; averaging writes
it as $v-hv$, and the surjection $j_\Q$ lifts $v$ so that conjugating
the first splitting removes the discrepancy.

In compatible split coordinates a minimal automorphism acts by the
scalar $\mu$ on $M$, with the finite-component action determined by
$\nu_2$.  This finite action has an actual lift to $C$:
$\rho\mapsto\rho$, $\sigma\mapsto\omega\sigma$.
It preserves $K$, fixes $K^{\mathrm{ab}}$, and thus induces an
automorphism of $\mathscr C_2$.  After composing with this lift,
the remaining component action is the identity.  Its translation
part in split coordinates is a one-cocycle and hence a coboundary
$v-hv$.  Lift $v$ through $j_\Q$, and compose the resulting inner
automorphism with scalar multiplication by $\mu$ on $K^{\mathrm{ab}}$.
This is the required automorphism.  All averages and lifts can be
chosen rationally once and for all, so the construction works after
base change to every $\Q$-algebra.  It is an existence assertion;
different lift choices differ by automorphisms over the coarser model.

The last assertions follow directly from the cocycle construction or
the universal property of relative completion: the image of every
source word and the compatible quotient map are specified before
completion.  There is no reverse universal property that lifts an
arbitrary automorphism of $\mathscr C_0$.  B.24 gives an explicit
failure of precisely that reverse assertion.
\end{proof}

\begin{statement}{B.24d}{Open}{The retained conjugator images on the complete relative tower}
B.24 refutes an unrestricted integral action obtained by lifting the
minimal \S51 automorphism functor.  B.24a--c construct the first integral
relative stage, its characteristic-two liftable locus and lift fibres,
including all additive-polynomial infinitesimals in the marked kernel.
They do not identify that minimal input with every version of
$M_{\rho,\sigma}$ or $N_{\rho,\sigma}$ in \S49.

B.31 now specifies the actual effective signed source and all the
retained conjugator fibre products, proves preservation of every
Fermat kernel, and constructs their continuous relative actions
over every $\mathbb Q_\ell$, on the profinite relative tower, and
on the abelian integral stages over $\mathbb Z_\ell$.
B.32 identifies the entire abelian Fermat tower with $C/F''$ and
the relative pro-prime family with its precise finite-quotient class.
B.33 computes the full actual-source image on the first modular
stage as $\widehat\mathbb Z\rtimes\widehat\mathbb Z^\times$.
B.34 refutes descent of the entire profinite action to the rational
base, already on that first Lie layer.

The remaining comparison is to compute the images of the retained
normalizer and conjugator data on the $q_2$ stage and then on all
higher relative stages, and to classify the full action kernels
beyond the explicit word-congruence criteria. Compatibility of
the arbitrary additive-polynomial first-stage lifts in B.24b
with every higher nonabelian stage also remains unproved.
Neither the nonempty rational lift fibres of B.24c nor the
functorial existence of the actual local action answers those
image and realization questions.
\end{statement}
The obstruction in the previous Open version of B.24 is thus sharpened:
the unrestricted minimal integral map does not exist, rather than
merely lacking a recovered infinitesimal parameter.  The explicit
liftable subfunctor and its computed first-stage kernel are the
replacement. The subsequent source-preservation, first-image, and
base-field results remove those former residuals; the retained
conjugator images and higher realization remain separate obligations.
